\begin{recipe}{Broodje hamburger}
\kicker[Hoofdgerecht,Vlees,Kinderen]{Hoofdgerecht · vlees · kinderen}
\index[register]{Broodje hamburger}
\index[register]{Hamburger}
\index[register]{Ui}

\meta{30 minuten}

\lettrine{B}{roodje} hamburger staat bij ons regelmatig op tafel. De winst zit in
kleine dingen: langzaam gekaramelliseerde uien, een geroosterd bolletje en een
plakje kaas dat op de hete burger smelt. De rest laten we door de kinderen zelf
samenstellen.

\blockrule{Ingrediënten}
\begin{ingredients}
  \ing{5}{hamburgers (black angus)}\index[register]{Hamburger}
  \ing{5}{hamburgerbolletjes}
  \ing{1}{ui}\index[register]{Ui}
  \ing{½ krop}{sla}
  \ing{1}{tomaat}
  \ing{3}{augurken, in plakjes}
  \ing{1}{komkommer, in plakjes}
  \ing{1}{paprika, in reepjes}
  \ing{5}{plakjes kaas}
  \ingb{mayonaise}
  \ingb{tomatenketchup}
  \ingb{burgersaus}
  \ingb{mosterd}
  \ingb{boter, om te bakken}
\end{ingredients}

\blockrule{Bereiding}
\tip{Snij de tomaat horizontaal door voor gelijkmatige, ronde plakken. Een vleestomaat
zoals coeur de bœuf geeft mooie grote plakken. Heb je wat meer tijd? Maak de
bolletjes dan zelf, met het recept voor hamburgerbolletjes (p.~\pageref{recipe:hamburgerbolletjes}).}
\begin{steps}
  \item Snij de ui in dunne ringen. Verhit een klontje boter met een scheutje water in
        een koekenpan op laag vuur en laat de uienringen langzaam karamelliseren. Het lage
        vuur en het scheutje water voorkomen dat de ui verbrandt voor hij zoet en
        zacht is.
  \item Verdeel de sla in blaadjes en snij de tomaat in plakken.
  \item Verhit bakboter in een tweede koekenpan en bak de hamburgers, ongeveer 10
        minuten. Check de verpakking voor de precieze tijd, dit verschilt per merk.
        Druk de burgers niet aan met de spatel: elk drukje perst het sap eruit dat
        in het broodje had moeten zitten.
  \item Rooster de bolletjes ondertussen onder de gril in de oven.
  \item Leg vlak voor de burgers klaar zijn een plakje kaas op elke hamburger. Doe
        voor het laatste beetje smelten kort een deksel op de pan.
  \item Laat de burgers 2 tot 3 minuten rusten (onder de kaas) voor je ze op het
        broodje legt: rust geldt ook voor klein vlees.
  \item Laat iedereen zijn eigen broodje opbouwen met wat hij of zij wil.
\end{steps}

\heroimagefade{images/broodje-hamburger}
\end{recipe}
